\section{Элементы теории графов}

\subsection{Основные понятия и определения}
Теория графов является важным инструментом для визуализации и анализа дискретных вероятностных пространств, марковских цепей и комбинаторных задач.

\begin{definition}
\textbf{Граф} $G = (V, E)$ — это математическая модель, состоящая из конечного непустого множества вершин $V$ и множества ребер $E$, представляющего собой связи между парами вершин. Формально, ребро — это неупорядоченная или упорядоченная пара вершин.
\end{definition}

\begin{itemize}
\item \textbf{Изолированная вершина} — вершина со степенью 0 (нет инцидентных ребер).
\item \textbf{Петля} — ребро, связывающее вершину саму с собой.
\item \textbf{Кратные ребра} — несколько ребер, соединяющих одну и ту же пару вершин.
\item \textbf{Простой граф} — граф без петель и кратных ребер. В курсе теории вероятностей чаще всего работают с простыми графами.
\end{itemize}

\subsection{Матричные способы задания графа}
Пусть граф имеет $n = \vert{}V\vert{}$ вершин и $m = \vert{}E\vert{}$ ребер.

\subsubsection{Матрица смежности $B$}
Размер матрицы $B$ составляет $n \times n$. Элемент $B_{ij}$ равен числу ребер, связывающих вершину $v_i$ и вершину $v_j$:
$$B_{ij} = \begin{cases} k, & \text{если вершины } v_i \text{ и } v_j \text{ соединены } k \text{ ребрами}, \\ 0, & \text{в противном случае.} \end{cases}$$
Для неориентированного графа матрица смежности всегда симметрична ($B_{ij} = B_{ji}$).

\subsubsection{Матрица инцидентности $A$}
Размер матрицы $A$ составляет $n \times m$. Для ориентированного графа (орграфа) элементы матрицы определяются следующим образом:
$$A_{ij} = \begin{cases} -1, & \text{если ребро } e_j \text{ выходит из вершины } v_i, \\ 1, & \text{если ребро } e_j \text{ входит в вершину } v_i, \\ 0, & \text{если ребро } e_j \text{ не инцидентно вершине } v_i. \end{cases}$$

\subsection{Связность, мосты и деревья}

\begin{definition}
Ребро $e$ связного графа называется \textbf{мостом}, если при его удалении граф распадается на несколько компонент связности.
\end{definition}

\begin{definition}
\textbf{Деревом} называется связный простой граф, не содержащий циклов. Деревья широко используются при построении деревьев вероятностей.
\end{definition}

\begin{theorem}[О числе ребер в дереве]
Любое дерево с $n \ge 1$ вершинами содержит ровно $m = n - 1$ ребер.
\end{theorem}

\begin{proof}
Доказательство проведем методом математической индукции по числу вершин $n$.
\begin{itemize}
\item \textbf{База индукции:} При $n = 1$ граф состоит из одной вершины. Циклов нет, ребер тоже нет. Формула дает $m = 1 - 1 = 0$. Утверждение верно.
\item \textbf{Индукционный переход:} Предположим, что любое дерево с $k$ вершинами имеет $k - 1$ ребер. Рассмотрим произвольное дерево $T$ с $k + 1$ вершинами. Известно, что в любом дереве с числом вершин больше одной есть хотя бы один «лист» (вершина степени 1).Удалим этот лист и инцидентное ему ребро. Полученный граф $T'$ останется связным (так как удаление листа не разрушает пути между оставшимися вершинами) и в нем не появится циклов (так как мы только удаляли элементы). Следовательно, $T'$ — тоже дерево.

По предположению индукции, так как в $T'$ ровно $k$ вершин, то в нем $k - 1$ ребер. Возвращая удаленную вершину и 1 ребро обратно, получаем, что в исходном дереве $T$ ровно $(k - 1) + 1 = k$ ребер.
\end{itemize}
Утверждение доказано для любого числа вершин $n$.
\end{proof}

\begin{definition}
\textbf{Полный граф} $K_n$ — простой граф с $n$ вершинами, в котором каждые две вершины соединены ребром. Число ребер в $K_n$ равно $m = \frac{n(n-1)}{2}$.
\end{definition}

\subsection{Планарные графы и формула Эйлера}

\begin{definition}
Граф называется \textbf{планарным}, если существует такое его изображение на плоскости, при котором ребра пересекаются только в своих концах (вершинах).
\end{definition}

\begin{theorem}[Формула Эйлера для плоских графов]
Для любого связного планарного графа, уложенного на плоскости, справедливо строгое соотношение:
$$V - E + F = 2,$$
где $V$ — число вершин, $E$ — число ребер, $F$ — число граней (включая одну внешнюю бесконечную грань).
\end{theorem}

\begin{proof}
Доказательство проведем индукцией по числу циклов в графе.
\begin{itemize}
\item \textbf{Случай 1 (без циклов):} Если граф является деревом, то по предыдущей теореме $E = V - 1$. Поскольку циклов нет, граф не ограничивает ни одной внутренней области, следовательно, существует только одна грань — внешняя ($F = 1$).
Проверяем формулу: $V - (V - 1) + 1 = V - V + 1 + 1 = 2$. Верно.
\item \textbf{Случай 2 (с циклами):} Предположим, формула верна для графов с $c$ циклов. Пусть наш граф имеет $c + 1$ циклов. 
Удалим одно любое ребро, принадлежащее какому-либо циклу. 
При этом связность графа не нарушится (так как разорван лишь один обходной путь), но число ребер уменьшится на единицу: $E' = E - 1$. 
Удаление ребра из цикла объединяет две смежные грани (внутреннюю и внешнюю по отношению к удаленному ребру) в одну, следовательно, число граней также уменьшится на 1: $F' = F - 1$. 
Число вершин остается неизменным: $V' = V$.

По индукционному предположению для нового графа выполняется: $V' - E' + F' = 2$.
Подставим исходные переменные: $V - (E - 1) + (F - 1) = V - E + 1 + F - 1 = V - E + F$.
Значение выражения не изменилось и по-прежнему равно 2.
\end{itemize}
Таким образом, формула доказана.
\end{proof}

\begin{theorem}[Теорема Понтрягина — Куратовского]
Граф является планарным тогда и только тогда, когда он не содержит подграфов, гомеоморфных полному графу $K_5$ или полному двудольному графу $K_{3,3}$.
\end{theorem}

\begin{figure}[htbp]
\centering
\begin{tikzpicture}[scale=0.9]
% Полный граф K_5
\begin{scope}[shift={(-4,0)}]
\node at (0,-2.3) {\textbf{а) Полный граф $K_5$}};
\foreach \i in {1,...,5} {
  \node[circle,fill=black,inner sep=2pt] (v\i) at (90-\i*72:1.5) {};
}
\foreach \i in {1,...,5} {
  \foreach \j in {\i,...,5} {
    \draw[thick,blue!70!black] (v\i) -- (v\j);
  }
}
\end{scope}

% Двудольный граф K_{3,3}
\begin{scope}[shift={(4,0)}]
\node at (0,-2.3) {\textbf{б) Двудольный граф $K_{3,3}$}};
\foreach \i in {1,2,3} {
  \node[circle,fill=black,inner sep=2pt] (a\i) at (-1.3, 2-\i*1) {};
  \node[circle,fill=black,inner sep=2pt] (b\i) at (1.3, 2-\i*1) {};
}
\foreach \i in {1,2,3} {
  \foreach \j in {1,2,3} {
    \draw[thick,red!70!black] (a\i) -- (b\j);
  }
}
\end{scope}
\end{tikzpicture}
\caption{Классические непланарные графы (критерий Куратовского)}
\end{figure}